% Einheit 3 von 4 – Parameter und Transformationen (bank/trigonometrische-funktionen/e3.jsonl, zusammenbau v0.1)
%% TODO Zweigzeile Teil 1: was hier gelernt wird (Satz in Schülersprache aus dem Einheitstitel) – Einheit 3
\einheitenkopf[e3]{Einheit 3 von 4 · Parameter und Transformationen}
\zweigzeile{ab Kl. 10 · keine P10-Aufgabe}
\setcounter{aufgabe}{17}

%% TODO Titel: Ich-kann-Satz fehlt in der Bank (Platzhalter: Kettenname) – Nr. 18
%% TODO Anweisung: ein Satz über der ersten Teilaufgabe fehlt in der Bank – Nr. 18
\begin{aufgabe}{Parameter und Transformationen}
\begin{teile}
\teil Höher oder schmaler? Was hat sich bei g gegenüber f geändert?\\ \kreuz{die Höhe (a)}\kreuz{die Breite (b)}
\end{teile}
%% TODO Darstellung für schwach fehlt in der Bank (trigonometrische-funktionen-e3-k1-s1-v1; Abschnitt „Für schwache Schüler“)
\swz{y = 4 · sin(x) – wie groß ist die Amplitude?}{}
%% TODO Darstellung für schwach fehlt in der Bank (trigonometrische-funktionen-e3-k1-s1-v2; Abschnitt „Für schwache Schüler“)
\swz{y = 2,5 · sin(x) – wie groß ist die Amplitude?}{}
%% TODO Darstellung für schwach fehlt in der Bank (trigonometrische-funktionen-e3-k1-s1-v3; Abschnitt „Für schwache Schüler“)
\swz{y = 6 · sin(x) – wie groß ist die Amplitude?}{}
%% TODO Darstellung für schwach fehlt in der Bank (trigonometrische-funktionen-e3-k1-s1-v4; Abschnitt „Für schwache Schüler“)
\swz{y = 1,5 · sin(x) – wie groß ist die Amplitude?}{}
%% TODO Darstellung für schwach fehlt in der Bank (trigonometrische-funktionen-e3-k1-s1-v5; Abschnitt „Für schwache Schüler“)
\swz{y = 8 · sin(x) – wie groß ist die Amplitude?}{}
\swfrage{Was bleibt gleich, was ändert sich?}
%% TODO Darstellung für schwach fehlt in der Bank (trigonometrische-funktionen-e3-k1-s2-v1; Abschnitt „Für schwache Schüler“)
\swz{y = 5 · sin(x) – zwischen welchen Werten schwingt die Welle?}{}
\swz{Lies die Amplitude am Graphen ab.}{\begin{ksys}[xmin=0,xmax=360,ymin=-3,ymax=3,xstep=45,ystep=0.5,ablesen]\funktion{2.5*sin(\x)}{f}\end{ksys}}
%% TODO Darstellung für schwach fehlt in der Bank (trigonometrische-funktionen-e3-k1-s4-v1; Abschnitt „Für schwache Schüler“)
\swz{y = −2 · sin(x) – was ist anders als bei y = 2 · sin(x)?}{}
%% TODO Darstellung für schwach fehlt in der Bank (trigonometrische-funktionen-e3-k1-s5-v1; Abschnitt „Für schwache Schüler“)
\swz{y = sin(3 · x) – wie lang ist die Periode?}{}
\swz{Lies die Periode ab und berechne b in y = sin(b · x).}{\begin{ksys}[xmin=0,xmax=360,ymin=-2,ymax=2,xstep=45,ystep=0.5,ablesen]\funktion{sin(4*\x)}{f}\end{ksys}}
\swa{Skizziere den Graphen von y = 2 · sin(3 · x) im ersten Vollkreis ab 0°.}{\begin{ksys}[xmin=0,xmax=360,ymin=-3,ymax=3,xstep=30,ystep=1]\end{ksys}}
\swz{Stelle die Gleichung y = a · sin(b · x) zum Graphen auf: erst Amplitude, dann Periode, dann b.}{\begin{ksys}[xmin=0,xmax=480,ymin=-5,ymax=5,xstep=60,ystep=1,ablesen]\funktion{4*sin(1.5*\x)}{f}\end{ksys}}
\begin{teile}
\teil Welche Gleichung gehört zu g?\\ \kreuz{y = 2 · sin(x)}\\ \kreuz{y = sin(4 · x)}\\ \kreuz{y = 4 · sin(x)}
\end{teile}
%% TODO Darstellung für schwach fehlt in der Bank (trigonometrische-funktionen-e3-k1-s10-v1; Abschnitt „Für schwache Schüler“)
\swz[3]{Der Faktor vor sin steigt von 1 auf 4. Beschreibe, was mit der Welle passiert.}{}
%% TODO Darstellung für schwach fehlt in der Bank (trigonometrische-funktionen-e3-k1-s11-v1; Abschnitt „Für schwache Schüler“)
\swz{y = sin(x) + 2 – wohin ist die Welle verschoben, zwischen welchen Werten schwingt sie?}{}
%% TODO Darstellung für schwach fehlt in der Bank (trigonometrische-funktionen-e3-k1-s12-v1; Abschnitt „Für schwache Schüler“)
\swz{Um wie viel Grad muss man y = sin(x) nach links schieben, damit y = cos(x) entsteht? Schreibe die Gleichung mit c.}{}
%% TODO Darstellung für schwach fehlt in der Bank (trigonometrische-funktionen-e3-k1-s13-v1; Abschnitt „Für schwache Schüler“)
\swz{y = 2 · sin(4 · x) – erste Nullstelle nach 0° und erste Hochstelle?}{}
\swa{a) Stelle zum Graphen f die Gleichung y = a · sin(b · x) auf. b) Skizziere daneben den Graphen von y = 2 · sin(1,5 · x).}{\begin{ksys}[xmin=0,xmax=360,ymin=-4,ymax=4,xstep=45,ystep=1,ablesen]\funktion{3.5*sin(3*\x)}{f}\end{ksys} \begin{ksys}[xmin=0,xmax=480,ymin=-3,ymax=3,xstep=60,ystep=1]\end{ksys}}
\end{aufgabe}

%% TODO Titel: Ich-kann-Satz fehlt in der Bank (Platzhalter: Kettenname) – Nr. 19
%% TODO Anweisung: ein Satz über der ersten Teilaufgabe fehlt in der Bank – Nr. 19
\begin{aufgabe}{zwei Graphen mit verschiedenen Parametern vergleichen}
\swz[3]{Vergleiche f und g im Bild: Was ist gleich, was ist verschieden?}{\begin{ksys}[xmin=0,xmax=360,ymin=-3,ymax=3,xstep=45,ystep=1,ablesen]\funktion{2*sin(\x)}{f}\funktion{2*sin(3*\x)}{g}\end{ksys}}
\end{aufgabe}

%% TODO Titel: Ich-kann-Satz fehlt in der Bank (Platzhalter: Kettenname) – Nr. 20
%% TODO Anweisung: ein Satz über der ersten Teilaufgabe fehlt in der Bank – Nr. 20
\begin{aufgabe}{Parameter und Transformationen – Fehler finden, Begründen}
\swz[3]{Lisa schreibt zu y = sin(4 · x): Die Periode ist 1440°, denn b = 4 macht die Welle breiter. Finde den Fehler.}{}
\swfrage{Begründe, warum ein größeres b die Welle schmaler macht.}
\end{aufgabe}

%% TODO Merkkasten Einheit 3: 7 Zeilen, Vorgabe höchstens fünf (3.1)
\uebersichtskasten{Allgemeine Sinusfunktion: y = a · sin(b · x) – a verändert die Höhe der Welle, b ihre Breite. \\ Amplitude a: Der größte Ausschlag nach oben und nach unten. Ist a negativ, wird die Welle an der waagerechten Achse gespiegelt. \\ y = 3 · sin(x) schwingt zwischen −3 und 3. \\ Faktor b: Die Periode wird durch b geteilt – Periode = 360° : b. Ein großes b macht die Welle schmal, ein b unter eins macht sie breit. \\ y = sin(2 · x) hat die Periode 180°, es passen zwei ganze Wellen in einen Vollkreis. \\ Gleichung aus dem Graphen: erst die Amplitude an Hoch- und Tiefpunkt ablesen, dann die Periode von Hochpunkt zu Hochpunkt messen, daraus b berechnen. \\ Nur Gymnasium: c verschiebt die Welle zur Seite, d nach oben oder unten – y = a · sin(b · x + c) + d.}
